\scanpage{281}

\unsure{Je présume} que dans le cas général \unsure{où on ne} suppose pas $\Pi_1(\mathcal A^0)=0$, il faut remplacer la condition $\Pi_1(\mathcal A^{0u})=0$ par
\[
\Pi_1(\mathcal A^{0u})\longrightarrow\Pi_1(\mathcal A^0)
\]
un isomorphisme.

\margin{Trop brutal ! La condition en question disant $\Pi_1(\mathcal A^{0\Gamma})\xrightarrow{\sim}\Pi_{\unsure{1}}(\mathcal A^0)^\Gamma$ -- mais il faut reprendre avec soin l'ensemble des conditions et les modifier ad hoc\ldots{} Cf\ldots{} feuille \unsure{intercalaire|intercalée}}

(Le plus agréable serait que ce soit un \unsure{homotopisme} -- c'est \unsure{cela} sans doute qui exprimerait qu'on a une $\infty$-\unsure{équivalence}\ldots)

J'ai envie de prouver \inserted{d'abord} $1^\circ$, $2^\circ$, $3^\circ$ dans le cas $\mathcal A=\mathcal A_{g,\nu}$, en me \unsure{limitant} au besoin au cas \inserted{anabélien} (sans doute le plus dur en fait ! \unsure{\}} mais \deleted{plus} \inserted{moins} touffu conceptuellement), et bien sûr au cas où $\Gamma$ est fini. Le premier travail sera bien sûr \unsure{celui|alors} de déterminer $\mathcal A^{0u}$ et sa structure topologique, pour essayer de prouver que $\mathcal A^{0u}$ est contractile dans ce cas

\scanpage{282}
\mspage{142}

\uuline{24} \inserted{Essai de} Détermination \unsure{de|des} $\mathcal A^{0\Gamma}$. Lien avec la relation $\Pi_{g,(\nu,\nu+n-1)}^\Gamma=(1)$, et programme de travail\ldots

Considérons une opération (fidèle \unsure{si on y tient}) du groupe fini $\Gamma$ sur la surface $\unsure{U}\subset X=\widehat U$ de type $(g,\nu)$, \inserted{$U=X-S$,} soit $\mathcal A=\operatorname{Aut}(\unsure{U})$, donc $\Gamma\subset\mathcal A$ et $\mathcal A^{u}=\mathcal A^\Gamma$ est le centralisateur de $\Gamma$. \inserted{Supposons d'abord $\Gamma=\Gamma^+$.} Evidemment $f\in\mathcal A^\Gamma$ implique que $f$ \deleted{\unsure{transforme} $\unsure{U}$ en lui $\overline{\mathrm m}$} \inserted{induit un \unsure{autom.} $g$ de $X/\Gamma\unsure{\equiv|\simeq}X'$} \deleted{\unsure{étant}} $X$ est un rev. galoisien de groupe $\Gamma$.

Soit $S'$ l'image de $S$ dans $X'$ (donc $S=X\mid S'$), et soit $S^!$ l'ens des pts de ramification de $\Gamma$ dans $U$, $S^{\prime!}$ son image dans $U'=X'\backslash S'$ (NB $U=X\mid U'$). Ainsi $g$ est un autom. de $(X',S')$, appliquant $S^{\prime!}$ en lui $\overline{\mathrm m}$, et \deleted{\unsure{tel que}} induisant donc un autom de $V'=U'\backslash S^{\prime!}=X'-(S'\cup S^{\prime!})$, \deleted{\unsure{sur lequel}} \inserted{L'image} \deleted{\illegible{2}} \inserted{inverse} $V=X\mid V'$ est étale sur $V'$, c'est $\overline{\mathrm m}$ un $\Gamma_{V'}$-torseur sur $V'$ et l'image inverse de celui-ci par $g\mid V'$ est (via $f$) isomorphe à $V$. La donnée de $f$ équivaut à la donnée d'un autom du \deleted{\illegible{1}} $\Gamma$-rev. $V$ de $V'$, induisant un

\scanpage{283}

isom de $V'$ qui, prolongé à $X'$, envoie $S'$ dans lui-$\overline{\mathrm m}$. Cela donne une interprétation de $\mathcal A^\Gamma$ en termes de constructions sur $X'$.

Considérons le groupe fond. $\Pi'$ de $V'$, \unsure{avec la} structure à lacets indexée par $S'\amalg S^{\prime!}$, on a donc un hom surjectif
\[
\Pi'\xrightarrow{\varphi}\Gamma
\]
qui sur \inserted{chacun} des ss-groupes \unsure{à|de} lacets d'indice $s'\in S^{\prime!}$ n'est pas triviale, et un autom $g$ de $V$ définit un autom extérieur \inserted{$\unsure{\mathfrak g|\overline g}$} de $\Pi'$ -- la condition à mettre dessus est que son composé avec $\varphi$ est conjugué \unsure{à|de} $\varphi$, et que $\unsure{\mathfrak g|\overline g}$ applique $S^{\prime!}$ dans lui-$\overline{\mathrm m}$.

\margin{on choisit un pt $a\in V$ et son image $a'\in V'$ comme pts bases, pour expliciter \illegible{1} $\unsure{\pi|\Pi'}$}

Ceci posé, \deleted{$f$ comme \unsure{$g$}} $g$ est défini par $\unsure{\mathfrak g|\overline g}$ à isotopie près \deleted{\unsure{soit|et}} (si on excepte des cas \inserted{$(g',\nu')$} dégénérés \deleted{\illegible{2}}, savoir $(0,\nu')$ avec $\nu'=0,1,2$ qu'il faudra examiner à part\ldots) et $g$ étant choisi, $f$ est défini modulo multiplication par un élément du centre de $\Gamma$, mais on sait déjà que $\mathcal A^0\cap\Gamma=(1)$, donc si on se borne à examiner des élts $f\in\mathcal A^{0\Gamma}$ alors $f$ sera déterminé par $g$ de façon unique

\scanpage{284}
\mspage{143}

On est ainsi amené au Pb suivant :

Soit $V'=X'\setminus(S'\amalg S^{\prime!})$ une surface
\deleted{\unsure{orientable}} \inserted{admissible} de type $(g',\nu')$
($X'$ compacte, $\nu'=\operatorname{card} S'+\operatorname{card} S^{\prime!}$),
munie d'un \unsure{sous-ens.} $S^{\prime!}$ de l'ens de ses pts à l'infini
et d'un pt base $x'$, d'où $\Pi'=\Pi_1(\unsure{X'|V'},x')$.
On a donné un rev. galoisien \inserted{\unsure{connexe|nommé} $V$} de $V'$ de groupe
fini $\Gamma$, ponctué au dessus de $x'$, donc défini par un hom surjectif
\[
\Pi'\xrightarrow{\varphi}\Gamma
\]
et on suppose \unsure{$V'$|$V$} ramifié en chacun des $s'\in S^{\prime!}$,
i.e. (pour un choix du \unsure{groupe à lacets|groupe lacets} \deleted{\illegible{1}} $L_{s'}$
correspondant à $s'$) que $L_{s'}\to\Gamma$ n'est pas trivial.
On considère \unsure{l'ens|l'espace} compactifié \deleted{\unsure{pour|par}} $X$ de $V$,
et l'ens $S$ resp $S^!$ des pts de $X$ sur $S'$ resp $S^{\prime!}$.
Soit $g$ un autom. de \unsure{$V$|$V'$}, définissant un autom ext.
$g_\bullet$ de $\Pi'$, on suppose que
\[
\unsure{g^*(V'/V)}\overset{\text{\inserted{\unsure{$V$-isom}}}}{\simeq}\unsure{V'}
\quad\text{i.e. que }\varphi\circ g_\bullet\text{ est conjugué à }\varphi
\]
\margin{OPS que $g$ fixe $x'$, donc $g_\bullet$ est un autom bien défini de $\Pi'$.}
\unsure{(|\{}ce qui exprime que $g$ provient par passage au quotient d'un autom $f$ de $V$,
défini mod multiplication par un \inserted{\unsure{élément de}}
$z\in$ Centre $\Gamma$).
Soit $U=V\cup S^!=X\setminus S$ (\unsure{rev. ramifié} sur
$U'=V'\cup S^{\prime!}=X'\setminus S'$),
\inserted{\unsure{On veut}} \deleted{\illegible{1}} exprimer que
\deleted{\unsure{$g$ provient par un}} \inserted{parmi les $f$ qui remontent}
$g$, il y en a un (\unsure{néc.|nécess.} unique!) qui est \deleted{dans}
\unsure{isot.} à $1$ dans $\underline{\mathrm{Aut}}(U)$ --- i.e. qui induit
l'autom ext. trivial de $\Pi_1(U)$. On aimerait

\scanpage{285}

prouver que ce sont exactement ceux qui \deleted{\unsure{induisent}} \inserted{sont}
isotopes à $1$ dans $\underline{\mathrm{Aut}}(V')$ --- ou encore que pour un tel $f$,
$f$ est nécessairement isotope à $1$ dans $\underline{\mathrm{Aut}}(\unsure{V})$
\deleted{\unsure{ou $\underline{\mathrm{Aut}}(U)$}}
\inserted{i.e. est dans $\underline{\mathrm{Aut}}(\unsure{V})^{0}$
($\subset\underline{\mathrm{Aut}}(\unsure{V|U})^0$)}
\unsure{et} pas seulement dans $\underline{\mathrm{Aut}}(U)^0$.
Notons \inserted{lorsque $\underline{\mathrm{Aut}}(V')^0$ est $1$-connexe} que
$\underline{\mathrm{Aut}}(V')^0$ se relève (en vertu de \unsure{principes généraux}) en
\deleted{\unsure{$\underline{\mathrm{Aut}}$}}
\deleted{\unsure{$\underline{\mathrm{Aut}}$}}
\[
\underline{\mathrm{Aut}}(V')^0\longrightarrow
\underline{\mathrm{Aut}}(V)^{0\Gamma}\subset
\underline{\mathrm{Aut}}(U)^{0\Gamma}.
\]

Donc la condition énoncée sur $g$ d'isotopie à $1$ est certainement
\uline{suffisante}. Notons d'ailleurs que les résultats déjà obtenus impliquent que
\[
\underline{\mathrm{Aut}}(U)^{\deleted{\unsure{0}}\Gamma}
\cap\underline{\mathrm{Aut}}(U)^0
=\underline{\mathrm{Aut}}(U)^{0\Gamma}
\]
induit \uline{l'identité} sur $S^!$ ($=U^!$ \unsure{avec} une notation antérieure).
Or \deleted{le sous-gr.} soit $B$
\deleted{\unsure{$\subset\underline{\mathrm{Aut}}(U)^0$}}
\inserted{$\subset\underline{\mathrm{Aut}}(U)^0\cap
\underline{\mathrm{Aut}}(\unsure{U|V},S^!)_{\unsure{=|\epsilon}}$}
le sous-groupe des pts \unsure{qui} fixent les $s\in S^!$, de sorte que
\[
\underline{\mathrm{Aut}}(U)^0/B\simeq\underline{\mathrm{Mon}}(S^!,U),
\]
et $B^0=\underline{\mathrm{Aut}}(V)^0$ --- un argument connu nous montre que
\[
B/B^0\simeq\Pi_1(\underline{\mathrm{Mon}}(S^!,U)).
\]
\inserted{Comme $\underline{\mathrm{Aut}}(U)^{0\Gamma}\subset B$ i.e.
$\underline{\mathrm{Aut}}(U)^{0\Gamma}=B^\Gamma$, \unsure{on}}
\deleted{On a donc
\diagram{\begin{tikzcd}[ampersand replacement=\&,column sep=1.6em,row sep=1em]
B^{0\Gamma}
 \arrow[r,phantom,"\subset" description]
 \arrow[d,phantom,"{\rotatebox[origin=c]{90}{$\simeq$}}" description]
\& \underline{\mathrm{Aut}}(U)^{0\Gamma}
 \arrow[r,phantom,"=" description]
\& B^\Gamma
\\
\underline{\mathrm{Aut}}(V')^0 \& {} \& {}
\end{tikzcd}}
, \unsure{donc} ($B^{0\Gamma}$ \unsure{est|étant} \unsure{connexe} et \illegible{2} $B^\Gamma$)}
\deleted{$\underline{\mathrm{Aut}}$}
:
\diagram{\begin{tikzcd}[ampersand replacement=\&]
1 \arrow[r] \& \unsure{B^{0\Gamma^0}} \arrow[r]
  \arrow[d, phantom, "\rotatebox{90}{$\simeq$}" description]
\& B^\Gamma \arrow[r] \arrow[d, phantom, "\rotatebox{90}{$\simeq$}" description]
\& (B/B^0)^\Gamma \arrow[r]
\& H^1(\Gamma,B^0) \arrow[d, equal, "?"] \\
\& \underline{\mathrm{Aut}}(V)^0 \& \underline{\mathrm{Aut}}(U)^{0\Gamma} \&\& 1
\end{tikzcd}}
\margin{\unsure{déduit donc de}
\diagram{\begin{tikzcd}[ampersand replacement=\&]
1\arrow[r]\&B^0\arrow[r]\&B\arrow[r]\&B/B^0\arrow[r]\&1
\end{tikzcd}}}

\scanpage{286}
\mspage{144}

Donc on trouve que la composante neutre de
$\underline{\mathrm{Aut}}(U)^{0\Gamma}$ est $\simeq\underline{\mathrm{Aut}}(V)^0$
(\unsure{donc} est $\infty$-connexe si $V$ \unsure{anabélien}), et son $\Pi_0$ est
\[
\subset(B/B^0)^\Gamma\simeq\Pi_1(\mathrm{Mon}(S^!,U))^\Gamma.
\]
On voudrait prouver que le sous-groupe \unsure{de|du} \inserted{ce} $\Pi_1$ des
invariants sous $\Gamma$ est réduit à $1$,
\margin{\unsure{?}}
\deleted{\uline{et mieux que simultanément}}
\[
\uline{\unsure{H^0(\Gamma,\Pi_1(\mathrm{Mon}(S^!,U)))}}
\deleted{\uline{=H^1(\Gamma,\Pi_1(\quad))=1}}
\]

On aimerait que ceci soit vrai même indépendamment d'hypothèses de réalisabilité,
quand on a donné un hom de $\Gamma$ dans le groupe de Teichmüller d'un $V$,
et qu'on fait \unsure{des} hypothèses adéquates....

\uline{Conjecture} Soit $\Pi$ un groupe extérieur à lacets de type $(g,\nu)$,
$I$ l'ens \unsure{d'indices|d'indice} des \inserted{classes des} lacets,
$\Gamma$ un groupe fini opérant extérieurement sur $\Pi$.
Alors il existe un groupe extérieur à lacets $\Pi'$, \unsure{d'ens|d'un} $I'$ de classes
de lacets, une \unsure{opération|action} extérieure de $\Gamma$ sur $\Pi'$, une partie
$I^{\prime!}$ de $I'$ stable \unsure{par|pour} $\Gamma$, un homomorphisme \inserted{ext}
de "bouchage des trous de $I'$"
\[
\Pi'\longrightarrow\Pi
\]
compatible avec l'action de $\Gamma$, tels que

\scanpage{287}

$1^\circ)$ Le stabilisateur dans $\Gamma$ de $\forall$ élément de $I^{\prime!}$ soit non réduit à $1$

$2^\circ)$ L'extension de $\Gamma$ par $\Pi'$ déduite de l'action ext. de $\Gamma$ n'a pas d'éléments d'ordre fini $\ne 1$ (i.e. pour aucun \unsure{élém\textsuperscript{t}} de $\Gamma$ $\ne 1$, l'action extérieure \inserted{sur $\Pi'$} ne se réalise \unsure{en|par} un autom. de $\Pi'$ d'ordre fini)

De plus, le $\Gamma$-groupe ext. à lacets $\Pi'$, muni du morphisme $\Pi'\longrightarrow\Pi$, est déterminé à isom unique près.

\deleted{\uline{\unsure{Commentaire.}}} \uline{Commentaire.} L'existence de $\Pi'$, \deleted{et} $\Pi'\longrightarrow\Pi$ est évidente dans le cas ``réalisable''. L'unicité -- à isom unique près, \inserted{même} dans le cas réalisable n'est pas évidente, ni même prouvée. Le noyau de \deleted{\illegible{1}}
\diagram{\begin{tikzcd}[ampersand replacement=\&]
\mathrm{Autext}_{\mathrm{lac}}(\Pi',\Pi) \arrow[r] \arrow[d,equal] \& \mathrm{Autext}(\Pi) \\
\mathrm{Autext}_{\mathrm{lac}}(\Pi',I^{\prime!}) \&
\end{tikzcd}}
est justement le groupe $\Pi_1$ de tantôt
\margin{plutôt, une extension de \unsure{$\mathfrak{S}_{I'}$} par \unsure{ce} $\Pi_1$}
et pour montrer que le foncteur des \deleted{\illegible{1}} \inserted{\unsure{couples} $(\Pi',I^{\prime!})$ d'un} $\Gamma$-groupe \unsure{ext.|cat.} $\Pi'$, \deleted{\illegible{1}} et d'une partie

\scanpage{288}
\mspage{145}

$I^{\prime!}$ de $I(\Pi')$ stable par $\Gamma$, telle que l'extension correspondante de $\Gamma$ par $\Pi'$ soit ``sans torsion'' et que le stabilisateur de $\forall\ i'\in I^{\prime!}$ dans $\Gamma$ soit non trivial, vers les $\Gamma$-groupes extérieurs, soit (\unsure{non} seulement ess. surjectif, mais) pleinement fidèle, est déjà problématique. La fidélité signifie justement que $\Pi_1^\Gamma=1$, la pleine fidélité est plus \unsure{forte|fort} que le fait que \unsure{trivial}
\[
\mathrm{Autext}(\Pi',\Pi)^\Gamma\longrightarrow\mathrm{Autext}(\Pi)^\Gamma
\]
est surjectif (dans le cas réalisable, cette \inserted{surjectivité} serait conséquence du fait que $\forall$ \unsure{classe} d'homéo. commutant à $\Gamma$ contient un homéo commutant à $\Gamma$)
-- il faut ajouter à ceci que \unsure{toute|tout} autre relèvement de $\Gamma$ en \unsure{op.|gr.} ext. de $\Pi'$ [i.e. un $\Gamma\longrightarrow\mathrm{Autext}_{\mathrm{lac}}(\Pi',\Pi)$\inserted{\diagram{\begin{tikzcd}[ampersand replacement=\&,row sep=1.2em]
\mathrm{Autext}_{\mathrm{lac}}(\Pi',\Pi)
 \arrow[d,phantom,"\inserted{\unsure{\rotatebox[origin=c]{90}{$\simeq$}}}" description]
\\
\inserted{\mathrm{Autext}_{\mathrm{lac}}(\Pi',I)}
\end{tikzcd}}}], \deleted{\unsure{telle que}} \inserted{ayant la} \unsure{$=$|$\overline{\mathrm m}$} \uline{action de $\Gamma$} \uline{sur} $I^{\prime!}$, est conjugué \unsure{par|à} la \illegible{1}

\scanpage{289}
\unsure{précédente|précédent} par un élément du \unsure{noyau} $\Pi_1$ de
\[
\mathrm{Autext}_{\unsure{\mathrm{lac}|\mathrm{loc}}}(\Pi',I^{\prime!})
\longrightarrow
\mathrm{Autext}_{\unsure{\mathrm{lac}|\mathrm{loc}}}(\Pi)
\times \unsure{\mathfrak S}_{I^{\prime!}}\quad\ldots
\]

On va \unsure{réénoncer} la conjecture précédente sous une forme un peu plus générale.
On va introduire une notation. Si $\Gamma$ est un groupe fini opérant \unsure{extérieurement}
\inserted{de façon directe} sur un groupe à lacets $\Pi$ \inserted{\unsure{anabélien}},
donnant lieu à une extension $\unsure{\mathcal E|E}$ de $\Pi$ par $\Gamma$,
on \unsure{va} \unsure{désigner} par $\Phi=\Phi(\Pi,\Gamma)$ (``\unsure{pts} fixes'')
\unsure{l'ens} des classes de $\Pi$-conjugaison des sections partielles
\inserted{$\ne(1)$} maximales de l'extension.
(Le caractère intrinsèque de $\Phi(\Pi,\Gamma)$, indépendamment des choix particuliers de $\unsure{\mathcal E|E}$, a été déjà noté).
Il est clair que $\Phi$ est un $\Gamma$-ensemble, on sait aussi qu'il est fini.

Soit \unsure{maintenant} $(\Pi',\Gamma)$ un $\Gamma$-groupe \inserted{\unsure{ext}}
à lacets, \inserted{\unsure{avec} $\Gamma$ opérant de façon directe ($\Gamma=\Gamma^+$)
et \unsure{fidèle} (\unsure{pour simplifier})} \unsure{fixons-nous} une partie
$I^{\prime!}$ de $I'=I(\Pi')$, stable par $\Gamma$,
\inserted{telle que $\forall i'\in I^{\prime!}$, \unsure{on} ait $\Gamma_{i'}\ne(1)$,}
et considérons le groupe extérieur quotient, défini par bouchage de $I^{\prime!}$,
soit $\Pi$ \deleted{\unsure{$\Pi'_{I^{\prime!}}$}}, -- il est clair que $\Gamma$ opère encore dessus.
On veut d'abord définir une bijection

\scanpage{290}
\mspage{146}
\[
(*)\qquad\Phi(\Pi,\Gamma)\overset{\sim}{\longleftarrow}
\Phi(\Pi',\Gamma)\unsure{\amalg|\sqcup} I^{\prime!}
\]
\margin{en supposant $\Pi,\Gamma$ également
\unsure{anabélien}, pour être plus sûr !}

a) Application $I'\longrightarrow\Phi(\Pi,\Gamma)$

Soit $L_{i'}\subset\Pi'$ de \unsure{classe} $i'\in I'$, $Z_{i'}$ son centralisateur dans l'extension $\unsure{\mathcal E|E}'$ de $\Gamma$ par $\Pi'$
de sorte que $L_{i'}$ est une extension de $\Gamma_{i'}$
par $L_{i'}$. Le passage au quotient $\unsure{\mathcal E|E}'\longrightarrow\unsure{\mathcal E|E}$
définit une section partielle $\Gamma_{i'}\hookrightarrow\unsure{\mathcal E|E}$,
contenue dans une unique section partielle \unsure{maximale}
\inserted{(sans doute elle est déjà maximale elle-\unsure{$=$|$\overline{\mathrm m}$} ?\ldots)}.
\margin{Oui, car pour \illegible{2} \unsure{$\Gamma$ cyclique},
\illegible{5} \unsure{réalisable}\ldots ! \unsure{?}}
On trouve ainsi une application $I'\longrightarrow\Phi(\Pi,\Gamma)$,
évidemment compatible avec les actions de $\Gamma$.

b) Application $\Phi(\Pi',\Gamma)\longrightarrow\Phi(\Pi,\Gamma)$

Toute section partielle \inserted{$\ne1$} de $\unsure{\mathcal E|E}'$ sur $\Gamma$ en définit
évidemment une de $\unsure{\mathcal E|E}$ sur $\Gamma$ en passant au
quotient, on sait comme ci-dessus qu'une
section maximale donne une section maximale,
en se ramenant au cas $\Gamma$ cyclique.
\margin{pas clair, car \unsure{l'hyp.} que les $\Gamma_{i'}\ne(1)$
\unsure{risque|risquent} de ne pas être \unsure{réalisée|réalisé}}

c) Bijectivité de (*). \deleted{Pour la surjectivité,} \inserted{Elle n'est}
\deleted{on est \unsure{ramené} au cas}
\inserted{\unsure{prouvée} (\unsure{sauf} si la situation de départ est}
réalisable, ou si $I^{\prime!}$ réduit à un seul élément
\ldots mais alors $\Gamma$ cyclique et la situation est réalisable\ldots)
